% Generated from validated Article Draft Result. Do not hand-edit; revise the structured draft instead.
\section{ModelからActionへ — Browser・Research・Tool・Codingへ広がる実行面}
\noindent\textbf{2025年前半のAgent化は一つの完成形へ収束したのではなく、browser操作、調査、API orchestration、cloud codingという異なるaction surfaceが並行して立ち上がる過程だった。}\autocite{src-78593f4b387f,src-8fec3ef3a97d,src-95b210494c88}

1月23日のOperator research previewは、ブラウザ上の操作をuser-facing agentとして切り出した。会話応答ではなく画面上で行動するsurfaceが製品として前景化した点に意味がある。\autocite{src-95b210494c88}


2月2日のDeep Researchはmulti-step web researchをChatGPT内の仕事としてまとめた。Operatorがbrowser actionを前面に出したのに対し、Deep Researchは情報探索と統合を主目的とし、同じ「agent」でもaction semanticsは異なる。\autocite{src-8fec3ef3a97d}


3月11日のResponses API、built-in tools、Agents SDKは、web/file search、computer use、tracingなどをdeveloperが組み合わせるplatform primitivesとして提示した。Agentは単体アプリではなく、modelとtoolを束ねるharness設計の問題へ移った。\autocite{src-78593f4b387f}


5月にはOpenAI Codexのcloud software-engineering agentとGoogle Julesのpublic betaが並んだ。Codexはparallel sandbox tasks、JulesはGitHub repositoryをcloud VM上で扱う非同期作業として現れ、coding agentもCLI補助から遠隔実行環境へ広がった。\autocite{src-8b595b787bad,src-e7aef8f1b13b}


同じ5月、Responses APIはremote MCPやCode Interpreterなどを拡張し、AnthropicもMessages APIへweb search toolを追加した。外部情報源・tool surfaceへ接続するほど、agentの実用能力はbase model単体よりもexecution environmentの設計に依存する。\autocite{src-9257c39d4172,src-e31c642d34b6}


H1の各製品は互換な「Agent製品」ではない。browser action、research、orchestration、coding sandboxを区別したうえで見ると、半年の変化はmodel responseからaction surfaceへ能力の観測単位が広がったことにある。\autocite{src-78593f4b387f,src-8fec3ef3a97d,src-95b210494c88,src-e7aef8f1b13b}


\begin{claimboundary}
一次資料から確認できるのは、各製品・APIが提示したaction surfaceとavailabilityである。タスク成功率や生産性向上といった性能評価は、本章では独立検証済みの事実として扱わない。
\end{claimboundary}
